\section{A relative-kernel proof of periodicity}
\label{app:periodicity}

\begin{proof}[Proof of Lemma~\ref{lem:periodicity}]
Write $T_j=T_{\OO_X(j)}$ and $L=(-\otimes\OO_X(1))$.
We interpret the twist as the Fourier--Mukai transform whose kernel is
the cone of the evaluation map
\[
 \OO_X(-j)\boxtimes\OO_X(j)\longrightarrow\OO_\Delta.
\]
We take cones in the standard enhancement by complexes of perfect
kernels, so that the evaluation squares below induce maps of cones.
Their transforms give natural transformations of functors. Tensoring the evaluation
kernel gives $LT_jL^{-1}\simeq T_{j+1}$, and hence
\begin{equation}\label{eq:twist-product}
 \Phi^5\simeq T_0T_1T_2T_3T_4L^5.
\end{equation}
We will identify the kernel of $T_0T_1T_2T_3T_4$ with
$\OO_\Delta(-5)[2]$, where the twist is in the second factor.

Keep the first copy of $X$ as a source parameter. Set
\[
 f=1_X\times i\colon X\times X\longrightarrow X\times\mathbf P^4,
 \qquad
 p\colon X\times\mathbf P^4\longrightarrow X,
 \quad q\colon X\times X\longrightarrow X.
\]
Here $p$ and $q$ are the first projections.
The kernel convention is that $F\in\Db(\Coh(X\times Y))$ sends $E$ to
$R\pi_{Y*}(\pi_X^*E\otimes^{\mathbf L}F)$.
For such a kernel, let
\[
 a_j^Y(F)=R\pi_{X*}\bigl(F\otimes\pi_Y^*\OO_Y(-j)\bigr),
 \qquad
 \mathsf L_j^Y(F)=
 \operatorname{Cone}\bigl(a_j^Y(F)\boxtimes\OO_Y(j)\longrightarrow F\bigr).
\]
The arrow is the adjunction evaluation. Projection formula and proper
base change show that $\mathsf L_j^Y(F)$ represents postcomposition of
the transform of $F$ with the evaluation cone for $\OO_Y(j)$.
Indeed, on an input $E$ the first kernel in this cone gives
\[
 R\Gamma\bigl(X,E\otimes^{\mathbf L}a_j^Y(F)\bigr)\otimes\OO_Y(j)
 \simeq
 R\operatorname{Hom}\bigl(\OO_Y(j),\Phi_F(E)\bigr)\otimes\OO_Y(j).
\]
In particular the source coefficient here is $a_j^Y(F)$ itself.

These coefficient functors let us compare the evaluation cones on $X$ with
those on $\mathbf P^4$. We will show that the five projective-space
evaluations remove all coefficients and leave the zero kernel, while the
comparison preserves the cone of the divisor counit. Computing that cone
will give the required shift.

Start with
\[
 K_5=f_*\OO_\Delta,\qquad G_5=\OO_\Delta,\qquad
 u_5\colon Lf^*K_5\longrightarrow G_5
\]
given by the counit. For $j=4,3,\ldots,0$, form successively
\[
 K_j=\mathsf L_j^{\mathbf P^4}(K_{j+1}),
 \qquad
 G_j=\mathsf L_j^X(G_{j+1}).
\]
We construct compatible maps $u_j\colon Lf^*K_j\to G_j$ and show
that their cones are all isomorphic to $\operatorname{Cone}(u_5)$.

For any map $u\colon Lf^*K\to G$, the unit followed by $f_*u$ gives
$K\to f_*G$, hence comparison maps
\[
 \theta_l\colon a_l^{\mathbf P^4}(K)\longrightarrow a_l^X(G).
\]
For $(K_5,G_5,u_5)$ these are isomorphisms for every $l$: the composite
$f_*\OO_\Delta\to f_*Lf^*f_*\OO_\Delta\to f_*\OO_\Delta$
is the identity by the adjunction identity. Suppose that, before the
step indexed by $j$, the maps $\theta_l$ are isomorphisms for
$0\le l\le j$. The evaluation square
\[
\begin{array}{ccc}
 q^*a_j^{\mathbf P^4}(K_{j+1})\otimes\OO_X(j)
   &\longrightarrow& Lf^*K_{j+1}\\
 \big\downarrow\scriptstyle{\simeq}&&\big\downarrow\scriptstyle{u_{j+1}}\\
 q^*a_j^X(G_{j+1})\otimes\OO_X(j)
   &\longrightarrow& G_{j+1}
\end{array}
\]
commutes by adjunction. Taking its cones defines $u_j$; because its
left vertical arrow is an isomorphism, the cone of $u_j$ is isomorphic
to the cone of $u_{j+1}$.

For $l<j$, apply the coefficient functor $a_l$ to the two evaluation
triangles. The comparison on their first terms is
\[
 a_j^{\mathbf P^4}(K_{j+1})\otimes R\Gamma(\mathbf P^4,\OO(j-l))
 \longrightarrow
 a_j^X(G_{j+1})\otimes R\Gamma(X,\OO_X(j-l)).
\]
It is an isomorphism: $\theta_j$ is an isomorphism, and the divisor
sequence
\[
 0\longrightarrow\OO_{\mathbf P^4}(m-5)
 \longrightarrow\OO_{\mathbf P^4}(m)
 \longrightarrow i_*\OO_X(m)\longrightarrow0
\]
identifies the two cohomology complexes for $1\le m\le4$.
The comparison on the middle terms is also an isomorphism by induction,
so the remaining maps $\theta_l$ stay isomorphisms after the mutation.
This proves the induction and the constancy of the cone.

We next show that $K_0=0$. The step indexed by $j$ kills
$a_j^{\mathbf P^4}$, since
$R\Gamma(\mathbf P^4,\OO)=\C$. It preserves the already vanishing
coefficients with indices $k>j$, since
$R\Gamma(\mathbf P^4,\OO(j-k))=0$ for $1\le k-j\le4$.
Thus $a_j^{\mathbf P^4}(K_0)=0$ for $0\le j\le4$.
All these kernels are perfect, and proper base change implies that
every derived source fiber $F$ of $K_0$ satisfies
\[
 R\operatorname{Hom}_{\mathbf P^4}(\OO(j),F)=0
 \qquad(0\le j\le4).
\]
The exact Koszul complex of the five homogeneous coordinates propagates
these five consecutive vanishings to
$R\Gamma(\mathbf P^4,F(n))=0$ for every integer $n$.
For sufficiently large $n$, Serre vanishing and global generation,
applied to the finitely many cohomology sheaves of $F$, force all those
sheaves to vanish. Therefore every derived source fiber is zero, and
the derived Nakayama lemma gives $K_0=0$. It follows that
\[
 G_0\simeq\operatorname{Cone}(u_0)
       \simeq\operatorname{Cone}(u_5).
\]

Finally, $f$ is a Cartier divisor embedding with normal bundle
$\OO_X(5)$ in the target variable. Its two-term divisor resolution
computes
\[
 \mathcal H^{-1}(Lf^*f_*\OO_\Delta)=\OO_\Delta(-5),
 \qquad
 \mathcal H^0(Lf^*f_*\OO_\Delta)=\OO_\Delta,
\]
with no other cohomology sheaves. The counit $u_5$ induces the identity
in degree zero. Consequently its cone has the single cohomology sheaf
$\OO_\Delta(-5)$ in degree $-2$, and
\[
 G_0\simeq\OO_\Delta(-5)[2].
\]
Since $G_0$ is the composition kernel for $T_0T_1T_2T_3T_4$, this is
an isomorphism of kernels and therefore a natural functor isomorphism.
Equation~\eqref{eq:twist-product} now gives $\Phi^5\simeq[2]$.
The asserted inverse follows by composing this relation with $[-2]$.
\end{proof}
